\selectlanguage{english}
Grothendieck's appendix to Chapter XII examines a difficulty that becomes visible as soon as one tries to use small categories as models for spaces: an ordinary categorical pushout can discard the homotopy information that a homotopy pushout must retain. The investigation begins with a proposed contractibility criterion, passes through explicit finite counterexamples, and ends with a conjectural gluing theorem. Its exploratory order matters. An unsuccessful cohomology calculation and an impossible retraction are part of the route to the examples; correcting them changes some intermediate claims without destroying the main obstruction. The electronic edition already marks several of these revisions \citep{Grothendieck1990XIIapp}.

Throughout these pages \(BC=|NC|\) denotes the classifying space of a small category. The localizer \(\mathbf W_\infty\) consists of the functors inducing weak homotopy equivalences. The word ``aspherical'' means weakly contractible in this usage. A presheaf \(F:C^{\mathrm{op}}\to\Ens\) has a homotopy type through its category of elements \(C/F\). In particular, a representable presheaf is aspherical, since its category of elements is a slice \(C/c\) with a terminal object. These conventions connect the appendix to the integration of homotopy types described in \emph{Pursuing Stacks}, \S69, and to the adjoints-to-restriction viewpoint of the 1991 letter to Thomason \citep{Grothendieck1983,Grothendieck1991Letter}.

The opening hypothesis asks for two weakly terminal objects \(b,c\), meaning that every object admits a map to each, and for their categorical product \(a=b\times c\). The source proves that \(BC\) is simply connected. Its proof is sound: in the groupoid completion, cancellation against the two projections identifies all maps from an object to \(a\); those unique classes are compatible with every arrow and its inverse. The completion is therefore equivalent to the terminal groupoid. The proposed stronger conclusion, weak contractibility, is false, but the first example in the manuscript needs a different calculation from the one printed.

That example is the category \(X_G\) of the first two formal powers of an object with endomorphism monoid \(G\):
\[
 \Hom_{X_G}(b^m,b^n)=\bigl(\{1,\ldots,m\}\times G\bigr)^n,
 \qquad m,n\in\{1,2\}.
\]
It is an opfibration over the opposite of the full finite-set category on a singleton and a two-element set, with fibres the one-object categories of \(G\) and \(G^2\). The resulting cohomological spectral sequence is legitimate. The provisional derived-limit formula used to evaluate it is not: the transposition of the two-element set and the common retraction of its two inclusions supply relations that the proposed two-term complex has omitted. Even a constant coefficient diagram contradicts the displayed cokernel formula.

The commentary gives an explicit projective resolution for the actual indexing category. It proves that \(BX_G\) is \(2\)-connected for every monoid \(G\), while
\[
 H^3(BX_{C_2};\mathbb F_2)\cong\mathbb F_2.
\]
Thus the two-object, twenty-six-arrow category remains a finite counterexample to the opening proposition. Its obstruction occurs one degree later than the manuscript asserts. The repaired computation is supplied as a commentators' development, with complete proofs and a separate normalized-nerve verification. The standard spectral-sequence and homotopy-colimit ingredients are identified explicitly \citep{Cegarra2020,Thomason1979}.

A second development supplies the degree-two counterexample promised by footnote~2 under the exact product hypotheses. For an integer \(r\geq1\), adjoin a new source \(d\) to the category of the first two powers for the trivial monoid, with \(r\) maps \(d\to b\) and the corresponding \(r^2\) maps \(d\to b^2\). Composition selects coordinates. The product remains an actual categorical product, and its new source link retracts up to homotopy onto the poset of nonempty subsets of an \(r\)-element set of cardinality at most two. Its nerve is the subdivided complete graph. The resulting three-object category \(Y_r\) satisfies
\[
 BY_r\simeq\bigvee^{(r-1)(r-2)/2}S^2.
\]
For \(r=3\) this gives a twenty-one-arrow model of \(S^2\). It verifies the footnote without confusing the opening product condition with the weaker presheaf condition used subsequently.

The persistent indexing shape is the span \(\Psi=(b_0\leftarrow a_0\to c_0)\). For a span of representables in \(C^\wedge\), its value at \(x\) is the set of components of a bipartite graph: the vertices are the maps \(x\to b\) and \(x\to c\), and the edges are the maps \(x\to a\). A terminal ordinary pushout says that all these graphs are connected. It does not say that their cycles vanish. The commentary proves both useful consequences stated in the source: a terminal pushout forces simple connectivity of \(BC\), and if either arrow out of \(a\) is monic, the presheaf pushout is aspherical. The latter follows because the graph components are then stars, and Quillen's Theorem A applies to their incidence categories \citep{Quillen1973}.

The four-object example \(Z\) makes the missing cycles visible. It has two parallel arrows \(0\rightrightarrows a\), followed by arrows \(a\to b,c\) which equalize them. Its presheaf pushout \(b\amalg_a c\) is terminal. Its nerve, however, is two cones on the parallel-arrow circle, and hence \(BZ\cong S^2\). The proposed retraction onto the circle cannot be a functor: applying it to either equalizing relation would identify the two distinct arrows. Footnote 7 corrects the homotopy type, but the later statement \(\pi_1(Z)\cong\mathbb Z\) repeats the earlier error. We retain both printed readings and explain their mathematical status.

Adding two incomparable cone vertices to \(Z\) gives the six-object category \(Q\) with \(BQ\cong S^3\). This is the reliable finite obstruction in the later pushout square: its three input categories are contractible, while its ordinary categorical pushout is not. The assertion is established directly from the nerve, without an unrestricted converse to a comma-category theorem. Such a converse would be false: a connected acyclic noncontractible link can have contractible suspension. The appendix's extension to arbitrary \(W\subseteq W_1\) therefore does not follow from this example; the obstruction for this pushout lies in degree three.

Grothendieck then inserts this local defect into the boundary of a product \(\Phi_n\times\Phi_m\) of alternating directed intervals. The boundary has classifying space \(S^1\), while the product is contractible. Its slices over its categorical closure are contractible, but four of them have no terminal object. Accordingly the boundary inclusion satisfies the local asphericity condition of Chapter XII and fails the stronger Dwyer neighbourhood-adjoint condition. Two surrounding closed Dwyer inclusions allow the six-object defect to survive the enlargement. The necessary pushout comparison is Thomason's; Cisinski's later correction distinguishes Dwyer maps from their retracts \citep{Thomason1980,Cisinski1999Dwyer}. The resulting inclusion fails to preserve weak equivalences under cobase change. In contemporary terminology it is not an \(h\)-cofibration.

The same construction extends, with proof here, to the boundary of every product of at least two such intervals: fixing interior sinks in the remaining coordinates carries the same \(S^3\) defect into higher dimensions. Later theory gives two complementary perspectives. Pushouts along Dwyer maps are preserved even as \((\infty,1)\)-categorical pushouts \citep{HackneyOzornovaRiehlRovelli2024}; no model structure on \(\Cat\) with Thomason weak equivalences can be cartesian \citep[Proposition~4.3]{Pavlov2024Thomason}. These results explain the continuing significance of the appendix's distinction between contractible links and actual adjoint neighbourhoods.

The question whether total asphericity forces a terminal object has a simpler answer. Here ``total'' quantifies over every left Kan extension \(f_!1_A\), a convention expressly distinguished in the edition from the total asphericity of \emph{Pursuing Stacks}. Filtered categories satisfy this condition, because the category of elements of a filtered colimit of representables is filtered. Grothendieck already records the correction elsewhere in the manuscript \citep[Chapter II, Corollaries~2.17--2.18]{Grothendieck1990II}. Both \((\mathbb N,\leq)\) and the one-object category of an unsplit idempotent give examples without a terminal object. The surrounding classification is discussed in the commentary on Chapter~II, manuscript pages~35--36, and is not repeated here.

The closing conjecture concerns a different, homotopically correct gluing. If \(A\) is a full sieve of \(X\), \(B\) its complementary full cosieve, and \(E\) the category of arrows from \(A\) to \(B\) with ordinary commuting squares, then
\[
 BX\simeq\operatorname{hocolim}\bigl(BA\longleftarrow BE\longrightarrow BB\bigr).
\]
The comparison is defined by the natural transformation supplied by the crossing arrows. We prove the formula using mixed simplices of the nerve: each crosses from \(A\) to \(B\) once, and its barycentric coordinates split into a point on each side and an interval coordinate. Thomason's comparison identifies the resulting bisimplicial bridge with \(BE\). This settles the \(\mathbf W_\infty\) conjecture and supplies the whole hierarchy of local-system descent statements at once. The exact specialization is worked out here; the general comparison theorem predates the manuscript, and no claim about priority for this formulation is made.

Two qualifications are mathematically necessary. First, for \(\mathbf W_\infty\), the true gluing formula does not imply the subsidiary assertion that one of its bridge projections must be a weak equivalence whenever a boundary inclusion is. The acyclic-link example disproves that implication even for finite posets with contractible pieces and union. A stronger reversible localizer axiom would supply the implication, but \(\mathbf W_\infty\) does not satisfy that axiom. The numbered references in the appendix require caution: Chapters XV and XVI use different Loc~\((6)\) conventions, and the positions of the primes in their variants differ. The detailed commentary records the checked formulations without certifying which convention the appendix intended \citep{Grothendieck1990XV,GrothendieckDerivateursXVI}. Second, replacing crossing arrows by all paths in the completed space changes the bridge: for the parallel-arrow circle, two crossing arrows give \(S^1\), whereas an all-paths bridge gives a countable wedge of circles. The source's covering-descent argument is fully justified by gluing locally constant set-valued functors across commuting squares; higher local systems give its contemporary extension \citep{Lurie2009HTT}. The explicit directed attaching data are what make the final formula work.

\medskip
\noindent\emph{Commentary, audit, and mathematical developments AI-produced by GPT-6 Astra (OpenAI, ChatGPT Pro mode), checked and integrated by Claude (Anthropic). The proofs marked as commentators' developments assert no priority.}
